\section{B. Frobenius eigenvalues and the integral norm-one lattice}

We use the endomorphism of cohomology induced by the $q$-power
Frobenius morphism; its eigenvalues on $H^1$ of an abelian variety have
complex absolute value $q^{1/2}$. Equivalently we use geometric Frobenius
on the usual Galois representation. Throughout this section $q=p^n$,
$n\geq1$, and $\ell\ne p$. We allow a smooth projective scheme to be
geometrically disconnected in degree zero, as the question explicitly
allows finite \'etale schemes. All positive-degree examples below are
geometrically connected.

\paragraph{B.1a. Realization in every weight (Known; derivation supplied).}
\label{B1a}
For every algebraic integer $\lambda$ of weight $i\geq0$, namely
$|\sigma(\lambda)|=q^{i/2}$ for every complex embedding $\sigma$, there
exists a smooth projective $X/\mathbb F_q$ such that $\lambda$ occurs on
$H^i(X_{\overline{\mathbb F}_q},\mathbb Q_\ell)$. For $i>0$ we may take
$X=A_\alpha^i$, where $\alpha^i=\lambda$ and $A_\alpha$ belongs to the
Honda--Tate isogeny class of $\alpha$. Thus the weak conjecture holds with
the explicit exponent $m=1$ in every weight; no class number or unit
index enters that exponent.

\begin{proof}
Let $i>0$ and choose any root $\alpha$ of $T^i-\lambda$. Integrality is
transitive: $\alpha$ is integral over $\mathbb Z[\lambda]$, which is
integral over $\mathbb Z$. Every embedding of $\mathbb Q(\alpha)$ into
$\mathbb C$ sends $\alpha$ to a number whose $i$th power is a conjugate
of $\lambda$. Hence every such image has absolute value $q^{1/2}$.
The hypothesis of the Honda--Tate theorem is therefore satisfied,
without any restriction on $p$, $n$, or $i$.
By \citet[Th\'eor\`eme~1(i), p.~96]{Tate1971}, there is a simple
$A_\alpha/\mathbb F_q$ with $\alpha$ among the roots of its Frobenius
polynomial. We may fix an embedding $\overline{\mathbb Q}\hookrightarrow
\overline{\mathbb Q}_\ell$ when naming the root. If $v\ne0$ is an
$\alpha$-eigenvector in $H^1(A_\alpha,\overline{\mathbb Q}_\ell)$,
the tensor $v^{\otimes i}$ is nonzero, belongs to the K\"unneth direct
summand
\[
 H^1(A_\alpha,\overline{\mathbb Q}_\ell)^{\otimes i}
 \ \subset\ H^i(A_\alpha^i,\overline{\mathbb Q}_\ell),
\]
and has eigenvalue $\alpha^i=\lambda$. Notice that the $i$ factors are
different factors of the product; no nonzero exterior power of a single
eigenvector is asserted.

If $i=0$, Kronecker's elementary argument gives that $\lambda$ is a root
of unity. For completeness, all powers of $\lambda$ have degree bounded
by $[\mathbb Q(\lambda):\mathbb Q]$, and their monic minimal-polynomial
coefficients are bounded by the corresponding binomial coefficients.
Only finitely many such integral polynomials exist. The set of all their
roots is finite, so two powers of $\lambda$ coincide. If its order is
$r$, take $X=\operatorname{Spec}\mathbb F_{q^r}$. Frobenius acts on its
$r$ geometric points as an $r$-cycle; on $H^0$ the eigenvalues are all
$r$th roots of unity.

This realization argument is already contained, in an equivalent
restriction-of-scalars form, in the proof of
\citet[Proposition~2.6]{Milne1994}. We have separated the unconditional
existence argument from that proposition's use of the Tate conjecture
to classify arbitrary simple motives.
\end{proof}

\paragraph{B.1b. Explicit construction data and the semisimple motive
(Proved, using Honda--Tate and numerical semisimplicity).}
\label{B1b}
The preceding prescription specifies $X$ by a uniquely determined
isogeny class and a product operation; every member of that class works.
There is also a terminating, although inefficient, construction by
equations. Compute the minimal polynomial $P_\alpha$ and the division
algebra index $d$ from the local invariants in
\citet[Th\'eor\`eme~1(ii)]{Tate1971}. The desired class has
Frobenius polynomial $P_\alpha^d$ and dimension
$g=d[\mathbb Q(\alpha):\mathbb Q]/2$. Enumerate projective equations
over $\mathbb F_q$ together with regular group-law and inverse-law
presentations, retaining smooth, geometrically connected, commutative
group varieties of dimension $g$. These tests reduce to ideal
membership, smoothness, and finite-type algebra over a finite field.
Every abelian variety has such a finite presentation. Choose
$\ell^a$ larger than twice every coefficient bound
$\binom{2g}{j}q^{j/2}$ for its Frobenius polynomial. The finite
\'etale group scheme $A[\ell^a]$ and its Frobenius action can be
computed by finite algebra over finite fields. Its characteristic
polynomial modulo $\ell^a$ has a unique integral lift within those
bounds. Stop at the first polynomial $P_\alpha^d$.
Honda--Tate proves termination. This gives an effective specification
for every input $\lambda$; the general enumeration has not been run.

The numerical motive $N=h^1(A_\alpha)^{\otimes i}$ is semisimple by
\citet[Theorem~1]{Jannsen1992}. It is homogeneous of weight $i$ and has
$\lambda$ as an intrinsic Frobenius eigenvalue. One does not need an
unproved realization functor on all numerical motives for this last
assertion. The K\"unneth projectors of an abelian variety are algebraic;
the kernel from homological to numerical endomorphisms is nilpotent by
\citet[Corollary~1]{Jannsen1992}. Frobenius acts semisimply on the
cohomology of $A_\alpha^i$. Its nonzero spectral idempotent belonging to
the minimal polynomial of $\lambda$ therefore cannot disappear in a
nilpotent ideal. Consequently that eigenvalue remains on $N$. This
proves the semisimple-motive form of the weak conjecture with $m=1$.

\paragraph{B.2a. The exact set and the first weights (Proved).}
\label{B2a}
The exact set of eigenvalues occurring in degree $i\geq0$, as $X$ varies,
is the set of Weil $q$-integers of weight $i$. Sufficiency is B.1a.
Necessity is \citet[Th\'eor\`eme~(1.6), p.~276]{Deligne1974}, applied
to each smooth projective $X$. In particular weight $2$ is obtained
from $\alpha^2=\lambda$ and $A_\alpha^2$, and weight $3$ from
$\alpha^3=\lambda$ and $A_\alpha^3$. Real numbers
$\lambda=\pm q^{i/2}$, all supersingular cases, and all Jacobi sums
which satisfy the stated Weil-integer condition fall under the same
proof. There is no omitted exception for $p=2$, $p=3$, or nonsquare $q$.
If geometric connectedness were imposed also in degree zero, the exact
degree-zero set would instead be $\{1\}$; that is not the convention
of the question.

\paragraph{B.3a. Products of weight-one numbers and the role of the field
(Proved; Lead B2 refuted as an obstruction to realization).}
\label{B3a}
For $i>0$, every such $\lambda$ is a product of exactly $i$ weight-one
Weil $q$-integers: take $\alpha,\ldots,\alpha$ from B.1a. The additional
root of unity may be $1$. For $i=0$, the Artin construction gives the
root of unity itself. Thus every eigenvalue in question occurs in an
effective tensor construction from abelian and Artin motives over the
original field $\mathbb F_q$.

There is no class-group obstruction to this assertion, because the
weight-one factors were not required to lie in the original number
field. Even if we insist that the factors lie in that field after
changing the finite field, an explicit power is immediate: for $i>0$,
$\lambda$ itself is a weight-one Weil $q^i$-integer, and $\lambda^i$
is the product of $i$ copies of it. This proves the proposed conclusion
of Lead B1 with $m=i$ without splitting any divisors; B.1a is stronger.
The unrestricted root argument also explains why the earlier study's
designation of the strong conjecture in weights above one as open is
incorrect.

\paragraph{B.3b. The $\mathbb Q(\sqrt{-23})$ test, including an equation
(Proved and computed).}
\label{B3b}
Put $L=\mathbb Q(\sqrt{-23})$, $\omega=(1+\sqrt{-23})/2$,
$\lambda=1+\omega$, and $q=2$. Then
\[
 \omega^2-\omega+6=0,\qquad
 \lambda^2-3\lambda+8=0,\qquad
 \lambda\bar\lambda=8.
\]
The prime $\mathfrak p=(2,\omega-1)$ satisfies
$(\lambda)=\mathfrak p^3$: the conjugate prime $(2,\omega)$ does not
divide $\lambda$, and its absolute norm is $8$. No algebraic integer
of $L$ has norm $2$, since
\[
 N(a+b\omega)=a^2+ab+6b^2
 =\left(a+\frac b2\right)^2+\frac{23b^2}{4}
\]
cannot equal $2$ for integers $a,b$. Hence the fixed-field
factorization into weight-one Weil $2$-integers is impossible, although
unrestricted realization is possible.
We can also verify $h_L=3$ without relying on a table of quadratic
forms. The Minkowski bound is $(2/\pi)\sqrt{23}<4$, so every ideal
class has a representative of norm $1$, $2$, or $3$. Both $2$ and
$3$ split, and $(\omega)=(2,\omega)(3,\omega)$ identifies the
class of one prime over $3$ with the inverse class of one prime
over $2$. The latter class is nontrivial by the norm-$2$
calculation and has cube one by $(\lambda)=\mathfrak p^3$.
The possible representative classes are exactly its three powers.

Here is an actual variety over $\mathbb F_2$. Let
$\mathbb F_8=\mathbb F_2[t]/(t^3+t+1)$ and let
\[
 E/\mathbb F_8:\quad y^2+xy=x^3+x^2+(t+1).
\]
Its projective closure is nonsingular: in the affine chart a singular
point would have $x=0$ and then $y=0$ from the two derivatives, which
contradicts $t+1\ne0$; the unique point at infinity is nonsingular.
Thus it is a smooth plane cubic with a rational origin. Exhaustive
enumeration gives five affine points and the point at infinity, so
$\#E(\mathbb F_8)=6$. The Frobenius polynomial is
$T^2-3T+8$. Set
$A=\operatorname{Res}_{\mathbb F_8/\mathbb F_2}E$, an abelian
threefold. Its $H^1$ is the induced representation, so its Frobenius
polynomial is $T^6-3T^3+8$; this is precisely the
restriction-of-scalars calculation of
\citet[\S3, Lemma~1]{Tate1971}.

In fact $X=A$ already suffices. On $H^3(A)=\bigwedge^3H^1(A)$ the
characteristic polynomial is
\[
 (T^2-3T+8)(T^6-24T^3+512)^3.
\]
To verify this identity without approximating roots, the supplied
program forms the $20\times20$ exterior-cube matrix of the integral
companion matrix of $T^6-3T^3+8$ and factors its characteristic
polynomial over $\mathbb Q$. This proves that $\lambda$ occurs on
$H^3(A_{\overline{\mathbb F}_2},\mathbb Q_\ell)$ for every $\ell\ne2$.
The same program enumerates the three reduced primitive forms of
discriminant $-23$, agreeing with $h_L=3$.

\subsection{B.4. The exact integral cokernel and corrected lifting problem}

For this subsection $L/\mathbb Q$ is a finite Galois CM field,
$K=L^{\langle\tau\rangle}$, and $S$ is the set of primes of $L$ over
$p$. Valuations are integral valuations $\operatorname{ord}_{\mathfrak s}$,
not absolute values and not normalized slopes. Write
\[
 \begin{split}
 E_{L,S}&=\mathcal O_L[1/p]^\times,\quad E_L=\mathcal O_L^\times,
 \quad E_K=\mathcal O_K^\times,\\
 M_0&=\{u\in E_{L,S}:u\bar u=1\},\quad
 D^-=\{(a_{\mathfrak s})\in\mathbb Z^S:
 a_{\tau\mathfrak s}=-a_{\mathfrak s}\},\\
 \Lambda&=\operatorname{div}(M_0)\subset D^-,\qquad
 C_{L,p}=D^-/\Lambda.
 \end{split}
\]
These groups are independent of $n$ in $q=p^n$.

\paragraph{B.4a. Exact cokernel, including its $2$-part (Proved).}
\label{B4a}
Let $K_D$ be the kernel of
$D^-\to\operatorname{Cl}(L)$, $D\mapsto[\mathcal O_L(D)]$,
where we regard $D$ as a fractional ideal. For $D\in K_D$ choose
$\alpha\in L^\times$ with $(\alpha)=D$ and set
\[
 \delta(D)=[\alpha\bar\alpha]\in E_K^+/N_{L/K}(E_L).
\]
Then the complete answer is
\[
 \boxed{\ C_{L,p}=D^-/\ker(\delta:K_D\to E_K^+/N E_L)\ }.
\]
In particular there is the exact sequence
\[
 0\longrightarrow
 \frac{N(E_{L,S})\cap E_K}{N(E_L)}
 \longrightarrow C_{L,p}
 \longrightarrow\operatorname{im}(D^-\to\operatorname{Cl}(L))
 \longrightarrow0.
\]
The left group is elementary $2$-torsion. The displayed presentation,
and not merely the orders of the end groups, specifies any nontrivial
extension between the two parts.

\begin{proof}
The kernel of $\operatorname{div}:M_0\to D^-$ is $\mu(L)$.
Indeed such an element is an ordinary unit and all its complex
conjugates have absolute value one; the Kronecker argument in B.1a
applies. If $D=(\alpha)$ is anti-invariant, then
$\alpha\bar\alpha$ is a totally positive unit of $K$. Replacing
$\alpha$ by $e\alpha$, $e\in E_L$, changes its norm by $N(e)$, so
$\delta$ is well defined and a homomorphism. Its value is trivial
exactly when $\alpha/e$ has norm one for some $e\in E_L$; that is,
exactly when $D\in\Lambda$.

The image of $\delta$ is precisely
$(N(E_{L,S})\cap E_K)/N(E_L)$. One inclusion follows by choosing the
generator of $D$. Conversely, an $S$-unit with unit norm has an
anti-invariant divisor, including zero valuation at every
$\tau$-fixed prime. This proves the reverse inclusion and the exact
sequence. Finally every square of a unit of $K$ is its norm regarded
as a unit of $L$, so the left group is killed by $2$.

If $h$ is the exponent of $\operatorname{Cl}(L)$, then $2hD\in
\Lambda$ for every $D\in D^-$: write $(\beta)=hD$ and use
$\beta/\bar\beta$, whose divisor is $2hD$ and whose norm is one.
Thus $C_{L,p}$ is finite and killed by $2h$.
Finiteness of the ideal class group is
\citet[Theorems~4.3--4.4]{MilneANT2020}; no conjecture on units or
class groups is being used.
\end{proof}

\paragraph{B.4b. A finite presentation suitable for computation (Proved).}
\label{B4b}
Choose one prime in each conjugate pair and identify $D^-=\mathbb Z^s$.
Let the columns of an integer matrix $H$ be a basis of the full-rank
class-relation lattice $K_D$. Choose generators $\alpha_j$ of its
principal ideals. In a basis of the free group $E_K^+$ write the norms
$\alpha_j\bar\alpha_j$ as the columns of a matrix $V$, and let the
columns of $W$ generate the norm lattice $N(E_L)$. Then
\[
 \Lambda=H\{x\in\mathbb Z^s: Vx\in W\mathbb Z^t\},\qquad
 C_{L,p}=\mathbb Z^s/\Lambda.
\]
This is an exact Smith-normal-form problem, including the entire
$2$-primary part. Its proof is the identity $\Lambda=\ker\delta$ in
B.4a written in bases. Choices of principal generators change $V$
by columns of the norm lattice and do not change the answer. If
$s=0$, both $D^-$ and $C_{L,p}$ are zero. For imaginary quadratic $L$
and split $p=\mathfrak p\bar{\mathfrak p}$ the unit-norm kernel is
zero, so
\[
 C_{L,p}\simeq
 \mathbb Z/\operatorname{ord}([\mathfrak p]^2)\mathbb Z.
\]
This verifies the corresponding portion of R7 but makes its omitted
$2$-part explicit in general. In B.3b it gives $C_{L,2}=\mathbb Z/3$.

\paragraph{B.4c. The purported unpowered canonical point is not defined
(Refuted; correction to R7).}
\label{B4c}
There is no inclusion-valued homomorphism
$M_0/\mu(L)\to L^\times$, $[u]\mapsto u$. Already $[1]=[-1]$ while
$1\ne-1$. Consequently the ``point of $D(M_0/\mathrm{torsion})$
whose evaluation on $M_0$ is the inclusion'' used in R7 and its cited
proof does not exist. This is a type error, independently of any
class-group issue.

More invariantly, the genuine tautological point is
$f_0\in D(M_0)(\mathbb Q)$, the Galois-equivariant map
$u\mapsto u$. It does not lie in $D(M_0)^0$, because its restriction
to $\mu(L)$ is not trivial. Thus it has no lift from a connected
norm-one torus into that identity component, for any $q$.
For a concrete counterexample with vanishing divisor cokernel take
$L=\mathbb Q(i)$ and $p=5$. With
$a=(2+i)/(2-i)$ we have
$M_0=\mu_4\times\langle a\rangle$, $C_{L,5}=0$; nevertheless
$f_0(i)=i\ne1$, while every point of $D(M_0)^0$ evaluates $i$ as $1$.
The assertion cannot be repaired merely by assuming that the divisor
isogeny is an isomorphism.

There is a second distinction. If $M=M(q,L)$ and $G=D(M)$, with
$\varepsilon$ corresponding to $q\in M$, then
\[
 X^*((G/j(\mathbb G_m))^0)=M_0/\mu(L),\qquad
 X^*((\ker\varepsilon)^0)=(M/\mathbb Z q)/\mathrm{torsion}.
\]
These lattices are naturally isogenous, but cannot be identified
without accounting for the possible factor $2$. On $M_0$ the natural
anti-divisor map $m\mapsto\operatorname{div}(m/\bar m)$ is
$2\operatorname{div}$, not $\operatorname{div}$.

\paragraph{B.4d. Canonical powered points and an exact lifting criterion
(Proved).}
\label{B4d}
Put $r=|\mu(L)|$ and $T=D(M_0/\mu(L))$. For every integer $k$ divisible
by $r$ there is a canonical point $g_k\in T(\mathbb Q)$ defined by
$[u](g_k)=u^k$. Divisibility by $r$ is necessary and sufficient for
this formula to be independent of the representative. Suppose
$\tau\mathfrak p\ne\mathfrak p$, let
$D_{\mathfrak p}\subset\operatorname{Gal}(L/\mathbb Q)$ be its
decomposition group, and put $L_0=L^{D_{\mathfrak p}}$,
$K_0=L_0\cap K$. Set
\[
 S_0=\operatorname{Res}_{K_0/\mathbb Q}
       \operatorname{Res}^{1}_{L_0/K_0}\mathbb G_m.
\]
The inclusion $\Lambda\subset D^-$ defines a $\mathbb Q$-isogeny
$\phi:S_0\to T$. For $z\in L_0^\times$ with $z\bar z=1$ and
$D=\sum n_{\mathfrak s}\mathfrak s\in D^-$, define
\[
 e_z(D)=\prod_{\mathfrak s\in S/\langle\tau\rangle}
              \sigma_{\mathfrak s}(z)^{n_{\mathfrak s}},
 \qquad\sigma_{\mathfrak s}\mathfrak p=\mathfrak s.
\]
The product is independent of orientations and of representatives of
the cosets $\sigma_{\mathfrak s}D_{\mathfrak p}$. Then the full set of
rational lifts of $g_k$ is
\[
 \boxed{\ \{z\in L_0^\times: z\bar z=1,\quad
 e_z(\operatorname{div}u_j)=u_j^k\ (1\leq j\leq s)\}\ },
\]
where $u_1,\ldots,u_s$ is any basis of $M_0/\mu(L)$ with chosen
representatives. When nonempty it is a right torsor under
$\ker\phi(\mathbb Q)$, acting by multiplication.

\begin{proof}
The character lattice of $S_0$ is $D^-$: the quotient of
$\mathbb Z^S$ by the sums of paired basis vectors is identified with
$D^-$ by sending the class of one basis vector to the difference of
the paired vectors. Galois descent gives
$S_0(\mathbb Q)=\{z\in L_0^\times:z\bar z=1\}$ and the displayed
evaluation formula. Replacing $\mathfrak s$ by $\tau\mathfrak s$
changes both the exponent and the value to their inverses, so the
product is unchanged. Because $z$ is fixed by $D_{\mathfrak p}$,
coset representatives do not matter. Equality of points of $T$ is
equality of their values on a basis of its character group. This
gives exactly the equations above and proves their sufficiency as
well as necessity. The usual kernel action is free and transitive
on a nonempty fibre.
\end{proof}

\paragraph{B.4e. A finite arithmetic test and explicit exponent (Proved).}
\label{B4e}
In B.4d every lift satisfies
\[
 \operatorname{div}_L(z)=k(\mathfrak p-\tau\mathfrak p).
\]
For a fixed $k$, first test whether an element of $L_0$ with that
divisor and norm one exists. Integral ideals, class relations, and
the unit norm quotient in B.4a give exactly this test; the valuations
must also be divisible by the ramification index of $L/L_0$.
If none exists, there is no lift. If $z_0$ exists, every candidate is
$z_0\eta$ with $\eta\in\mu(L_0)$. Test the finitely many candidates
in the equations of B.4d. This determines the answer, with no
remaining existence problem or omitted sign.

For a universal sufficient exponent, let $h$ be the exponent of
$\operatorname{Cl}(L)$ and $d=|D_{\mathfrak p}|$. Then
\[
 \boxed{\ k=2hdr\ }
\]
always works. In particular the full criterion is much stronger
than an ``if and only if the isogeny is an isomorphism'' claim.

\begin{proof}
For a putative lift, comparison of divisors in the equations on
$\Lambda$ gives a linear map $D^-\to\mathbb Z^{\{
\text{finite primes of }L\}}$ agreeing with $k$ times inclusion on
the finite-index sublattice $\Lambda$. The target has no torsion, so
the maps agree on all of $D^-$. Evaluating on
$\mathfrak p-\tau\mathfrak p$ gives the required divisor of $z$.
The ratio of two norm-one elements with that divisor is an ordinary
norm-one unit, hence a root of unity. This proves the finite test.

Choose $\beta\in L^\times$ with
$(\beta)=h(\mathfrak p-\tau\mathfrak p)$ and put
\[
 c=N_{L/L_0}(\beta/\bar\beta),\qquad z=c^r.
\]
Since $\tau$ is central, $c\bar c=1$, and
$\operatorname{div}_L(c)=2hd(\mathfrak p-\tau\mathfrak p)$.
For every $u\in M_0$, the ratio
$e_c(\operatorname{div}u)/u^{2hd}$ has zero divisor and norm one;
it is therefore in $\mu(L)$. Raising to the $r$th power makes it
one, giving $e_z(\operatorname{div}u)=u^{2hdr}$ exactly.
Replacing $\beta$ by another generator changes $c$ by an ordinary
norm-one unit, so $c^r$ is independent of that choice.
The construction works for $p=2$ and $p=3$ as well. If
$\tau\mathfrak p=\mathfrak p$, transitivity makes $D^-=0$;
then $M_0=\mu(L)$, $T=1$, and the first canonical powered point
$g_r$ is the identity and has its unique lift to the trivial torus.
\end{proof}

\paragraph{B.4f. Normalization inside $\ker\varepsilon$ and field extension
(Proved; the literal unpowered request is refuted).}
\label{B4f}
The natural map to $U^0=(\ker\varepsilon)^0$ is instead dual to
\[
 (M/\mathbb Zq)/\mathrm{torsion}\longrightarrow D^-,
 \qquad[m]\longmapsto\operatorname{div}(m/\bar m).
\]
It is injective with finite cokernel: its kernel before taking the
torsion-free quotient consists exactly of $m$ whose $m/\bar m$ is a
root of unity, and its image contains $2\Lambda$. For an integer
$k$ divisible by $r$, the rational normalized point
\[
 f^{2k}j(q^{-k})\in U^0(\mathbb Q)
\]
evaluates a character $m$ as $(m/\bar m)^k$. Its lifting criterion
is B.4d with $u_j$ replaced by generators of
$\{m/\bar m:m\in M\}$. The same construction of B.4e lifts it for
$k=2hdr$. This gives a fully typed replacement for the erroneous
unpowered normalization on 15:112 and in R7.

We stress the direction of base change. In the category obtained by
extending a fixed object from $\mathbb F_q$ to $\mathbb F_{q^k}$,
Frobenius is replaced by its $k$th power, and the character map is
$u\mapsto u^k$. The powered assertions above apply in exactly that
sense. Merely replacing the symbol $q$ by $q^k$ while retaining the
full norm-one group $M_0(q,L)=M_0(p,L)$ changes neither the lattice
nor its cokernel and does not create an unpowered canonical point.
For instance $C_{\mathbb Q(\sqrt{-23}),2}=\mathbb Z/3$ for every
$q=2^n$. This distinction is also why the normalized groups
$W_0^L(q)$ in \citet[Proposition~2.27]{Milne1994}, which impose an
additional integrality condition on normalized absolute values,
must not be substituted for the $M_0(q)$ of 15:110.

\paragraph{B.4g. Audit of the background results used here (Proved or
corrected as indicated).}
\label{B4g}
The principal-ideal construction in R5, from 15:96, is correct:
$\beta/\bar\beta$ has norm one and doubles an anti-divisor. We have
independently proved this exact step in B.4a. If no prime over $p$
splits in the CM extension, a norm-one $p$-unit has every valuation
zero and is a root of unity. For a Weil number of weight $i$ this
applied to $\lambda^2/q^i$ gives
$\lambda=q^{i/2}\zeta$ for a root of unity $\zeta$, in an algebraic
closure; no claim that the chosen square root lies in $L$ is needed.
The divisor and finite-index portions of R7 are verified by B.4a.
Its inclusion-valued point and the identification of two compact
tori are false as written, as B.4c proves. The proposed class-group
obstruction in Lead B2 is valid only as part of the fixed-lattice
calculation B.4a, not as an obstruction to B.2 or B.3.
